\begin{abstract}
A developer who sees a qubit end in the wrong state needs to know which lines of the program determine it,
the question program slicing answers for classical code. Slicing has been proposed for quantum programs,
but the dependence rules behind it have never been checked against what a quantum program computes, so a
developer cannot tell whether a slice shows the code that matters, and a slice can look right while hiding
the fault. We ask what a quantum slice should look like to support comprehension. A slice must preserve
the joint state of the criterion qubits for every input, and an operation is a dependence when changing it
can change that state. Under these definitions, the quantum dependence graph (QDG), in which every gate
reads and writes all of its qubits, gives sound and dependence-complete slices. It is the known backward
light cone, so what is new is the semantics that justifies it as a slice and the measurement of what the
alternatives cost. On \numNPairs{} circuit--qubit pairs, the classical program dependence graph, which reads
the control of a gate as unchanged, gives wrong slices for \numCpdgUnsoundPct{} of them, and
\numFdCpdgZMatchPct{} of these reproduce the measured outcomes, so a histogram check would accept them.
Entanglement edges, as in the quantum program dependence graph, keep \numRqOneYzSize{} of a circuit on
average, against \numRqOneLcSize{} for the QDG, only \numLcExcess{} percentage points above the dependences
themselves. Pruning by simulating one input shrinks slices by a further \numHTwoRelative{}, but its slices
hold only for that input and lose a fault that the QDG slice keeps in \numCQFour{} of \numRqThreeUnits{}
seeded-fault units. Slices narrow where a developer looks, but they do not improve mutation-based fault
ranking. Our static slicer for Qiskit, built on srcML, shows a slice as source lines, loop iterations and
qubit profiles, and we close with implications for comprehension tools and open problems.
\end{abstract}
